\documentclass[12pt]{article}

% page layout
\usepackage[a4paper, margin=1in]{geometry}
\usepackage{lineno}

% misc
\usepackage{multirow}
\usepackage{csquotes}

\usepackage{microtype}
% \overfullrule=5pt

% graphics and colors
\usepackage[dvipsnames]{xcolor}
\usepackage{graphicx}
\usepackage{float}

% small font size for captions
\usepackage[font=footnotesize,labelfont=bf]{caption}

% math and tables
\usepackage{amsmath}
\usepackage{booktabs}
\usepackage{siunitx}
\sisetup{retain-explicit-plus}
\DeclareSIUnit{\eggs}{eggs}
\DeclareSIUnit{\year}{year}
\DeclareSIUnit{\Otwo}{O\textsubscript{2}}

% bibliography
\usepackage[style=apa, backend=biber, uniquename=false, giveninits=false]{biblatex}
% \usepackage[style=nature, backend=biber, uniquename=false, giveninits=false]{biblatex}

% authors and affiliations
\usepackage{authblk}
\renewcommand*{\Affilfont}{\footnotesize\normalfont}

% listings and text code formatting packages
\usepackage{listings}

% hyperlinks
\usepackage{hyperref}
\hypersetup{
  colorlinks=true,
  linkcolor=Blue,
  filecolor=Maroon,
  citecolor=Blue,
  urlcolor=Blue
}

% custom commands for title page
\title{Beyond the average phenotype: individual variation alters mechanistic forecasts of climate responses}

\author[1*]{Urtzi Enriquez-Urzelai}
\author[1,2]{Gioele Pascucci}
\author[1]{Lumír Gvoždík}

\affil[1]{Czech Academy of Sciences, Institute of Vertebrate Biology, Brno, Czech Republic}
\affil[2]{Department of Botany and Zoology, Faculty of Science, Masaryk University, Brno, Czech Republic}
\affil[*]{Corresponding author: uenriquez@proton.me}
\date{\vspace{-5ex}}

\makeatletter
% fix authblk spacing
\renewcommand\AB@affilsepx{\protect\\\vspace{\baselineskip}\protect\Affilfont}

% store title and author for appendix
\let\papertitle\@title
\let\paperauthors\@author

% appendix title block
\newcommand{\appendixtitleblock}[1]{
  \begin{center}
    {\Large\textbf{Appendix #1} \textit{for:} \papertitle}
    \vskip 1em
      {\large\paperauthors}
  \end{center}
  \vskip 1.5em
}

% image width constraint
\def\maxwidth{%
  \ifdim\Gin@nat@width>\textwidth
    \textwidth
  \else
    \Gin@nat@width
  \fi
}
\makeatother

% custom counter for appendix cross-references
\newcounter{appnum}

% custom macro to setup the appendices
\newcommand{\appendixsetup}[1]{
  \setcounter{appnum}{\numexpr#1-1\relax}%
  \refstepcounter{appnum}
  \newpage
  \ifdefined\nolinenumbers\nolinenumbers\fi
  \setcounter{page}{1}
  \newrefsection
  \setcounter{section}{0}
  \renewcommand{\theHsection}{app#1.\arabic{section}}
  \renewcommand{\thefigure}{S#1.\arabic{figure}}
  \renewcommand{\theHfigure}{S#1.\arabic{figure}}
  \setcounter{figure}{0}
  \renewcommand{\thetable}{S#1.\arabic{table}}
  \renewcommand{\theHtable}{S#1.\arabic{table}}
  \setcounter{table}{0}
}
\addbibresource{references.bib}
\addbibresource{packages.bib}

\makeatletter
\def\maxwidth{ %
  \ifdim\Gin@nat@width>\linewidth
    \linewidth
  \else
    \Gin@nat@width
  \fi
}
\makeatother

\begin{document}







\appendixsetup{2}
\appendixtitleblock{2}\label{appendix-params}

{\centering \bfseries 
\Large DEB model parameters and example simulations
\par
}



\section{DEB core parameter analysis across individuals}

We estimated Dynamic Energy Budget (DEB) model parameters for 37 \emph{I.~alpestris} individuals using the standard DEB model (std) \parencite{kooijman2010}. The standard DEB model is governed by 12 primary parameters that jointly determine the rates of assimilation, mobilisation, maintenance, growth, and reproduction \parencite{marques2018}. For the population-average parameterisation, seven of the 12 primary parameters were freed for estimation; the remaining five were held at their literature values. At the individual level, four parameters were allowed to vary: the zoom factors $z$ (female) and $z_m$ (male), the energy conductance $\dot{v}$, and the volume-specific somatic maintenance rate $[\dot{p}_M]$. All remaining parameters were held at their population-average values.

The zoom factor $z$ scales maximum body size ($L_m = z \cdot L_{\text{ref}}$, with $L_{\text{ref}} = 1$~cm) and, through the physical covariation rules \parencite{kooijman2010}, simultaneously determines the maximum surface-specific assimilation rate $\{\dot{p}_{Am}\}$ and the maturity thresholds at birth ($E_H^b$) and puberty ($E_H^p$). Because $z$ (or $z_m$) and $[\dot{p}_M]$ vary among individuals while $\kappa$ is fixed, $\{\dot{p}_{Am}\}$ is not a free parameter but is derived from the individual estimates of $z$ and $[\dot{p}_M]$:

\begin{equation}\label{eq:zvpM-pAm}
\{\dot{p}_{Am}\} = \frac{z \cdot L_{\text{ref}} \cdot [\dot{p}_M]}{\kappa}
\qquad
  \mathrm{and}
\qquad
\{\dot{p}_{Am,m}\} = \frac{z_m \cdot L_{\text{ref}} \cdot [\dot{p}_M]}{\kappa}
\end{equation}

The complete set of parameter values is presented in \textbf{Table~\ref{tab:deb-params}}.

\begin{table}[H]
  \caption{DEB parameter values for \emph{I.~alpestris}: population-average estimates and individual ranges (min--max) across the 37 individuals.}
  \label{tab:deb-params}
  \centering\scriptsize\resizebox{\textwidth}{!}{%

\begin{tabular}[t]{lllrrr}
\toprule
\multicolumn{4}{c}{ } & \multicolumn{2}{c}{Individuals} \\
\cmidrule(l{3pt}r{3pt}){5-6}
\multirow{2}{*}[2.3em]{Parameter} & \multirow{2}{*}[2.3em]{Symbol} & \multirow{2}{*}[2.3em]{Units} & \multirow{2}{*}[2.3em]{Population average} & Min & Max\\
\midrule
\addlinespace[0.3em]
\multicolumn{6}{l}{\textbf{Primary parameters}}\\
\hspace{1em}max spec searching rate & $\{\dot{F}_m\}$ & l\,d$^{-1}$\,cm$^{-2}$ & 6.5 & $-$ & $-$\\
\hspace{1em}digestion efficiency of food to reserve & $\kappa_X$ & $-$ & 0.8 & $-$ & $-$\\
\hspace{1em}max surface-specific assimilation rate & $\{\dot{p}_{Am}\}$ & J\,d$^{-1}$\,cm$^{-2}$ & 265.7 & 217.6 & 303.3\\
\hspace{1em}energy conductance & $\dot{v}$ & cm\,d$^{-1}$ & 0.01914 & 0.01716 & 0.02134\\
\hspace{1em}allocation fraction to soma & $\kappa$ & $-$ & 0.5914 & $-$ & $-$\\
\hspace{1em}reproduction efficiency & $\kappa_R$ & $-$ & 0.95 & $-$ & $-$\\
\hspace{1em}vol-spec somatic maint & $[\dot{p}_M]$ & J\,d$^{-1}$\,cm$^{-3}$ & 138.4 & 107.4 & 162.8\\
\hspace{1em}surf-spec somatic maint & $\{\dot{p}_T\}$ & J\,d$^{-1}$\,cm$^{-2}$ & 0 & $-$ & $-$\\
\hspace{1em}maturity maint rate coefficient & $\dot{k}_J$ & d$^{-1}$ & 0.002 & $-$ & $-$\\
\hspace{1em}spec cost for structure & $[E_G]$ & J\,cm$^{-3}$ & 7302 & $-$ & $-$\\
\hspace{1em}maturity at birth & $E_H^b$ & J & 6.126 & $-$ & $-$\\
\hspace{1em}maturity at puberty for females & $E_H^p$ & J & $1.366 \times 10^{4}$ & $-$ & $-$\\
\addlinespace[0.3em]
\multicolumn{6}{l}{\textbf{Extended parameters}}\\
\hspace{1em}zoom factor & $z$ & $-$ & 1.135 & 1.102 & 1.198\\
\hspace{1em}zoom factor for males & $z_m$ & $-$ & 1.046 & 0.9978 & 1.118\\
\hspace{1em}max surface-specific assimilation rate & $\{\dot{p}_{Am,m}\}$ & J\,d$^{-1}$\,cm$^{-2}$ & 244.8 & 191.9 & 307.9\\
\hspace{1em}maturity at puberty for males & $E_{H,m}^p$ & J & 7474 & $-$ & $-$\\
\hspace{1em}maturity at metam & $E_H^j$ & J & 334.2 & $-$ & $-$\\
\hspace{1em}shape coefficient & $\delta_M$ & $-$ & 0.1878 & $-$ & $-$\\
\hspace{1em}Weibull aging acceleration & $\ddot{h}_a$ & d$^{-2}$ & $6.475 \times 10^{-9}$ & $-$ & $-$\\
\bottomrule
\end{tabular}

  }
\end{table}

\subsection{Parameter variability}



The coefficient of variation (CV) quantifies relative variability across individuals. \textbf{Figure~\ref{fig:zvpM-cv}} shows the CV values sorted in descending order. Among the four free parameters, $[\dot{p}_M]$ exhibited the largest relative variation across individuals, followed by $\dot{v}$, while $z$ was the most conservative. The mean individual values deviated from the population-average estimate by 0.9\% for \(z\), -0.9\% for \(z_m\), 1\% for \(\dot{v}\), and -0.8\% for \([\dot{p}_M]\).

\begin{figure}[H]
  \centering
  \includegraphics[width=0.6\textwidth]{/home/urtzi/Documents/papers/individuals_DEB/figures/1_base/param_zvpM_cv.pdf}
  \caption{Coefficient of variation for each varying parameter.
  $[\dot{p}_M]$ is the most variable parameter across individuals.}
  \label{fig:zvpM-cv}
\end{figure}

\textbf{Figure~\ref{fig:zvpM-distribution}} shows the distribution of each parameter with the population-average value marked for reference. Most parameters were broadly distributed around the population-average estimate, indicating that the individual-level data span a wider range of physiological strategies than the single population-average value suggests.

\begin{figure}[H]
  \centering
  \includegraphics[width=0.75\textwidth]{/home/urtzi/Documents/papers/individuals_DEB/figures/1_base/explore_zvpM_distribution.pdf}
  \caption{Distribution of each varying parameter. The dashed line marks the
  population-average estimate.}
  \label{fig:zvpM-distribution}
\end{figure}

\newpage 

\section{Fit of DEB models to the data}

In \textbf{Fig.~\ref{fig:fit-growth}} and \textbf{Fig.~\ref{fig:fit-metabolism}} we show the fit of the model with the final set of estimated parameters for the population-average parameterisation (for each sex in different colours) and for four selected females and males.


\begin{figure}[H]
  \begin{center}
    \includegraphics[width=\textwidth]{/home/urtzi/Documents/papers/individuals_DEB/figures/2_final/figS1-growth.pdf}
  \end{center}
  \caption{Model predictions of growth for the population-average parameterisation (\textit{blue}: males; \textit{red}: females) and eight selected individuals.}\label{fig:fit-growth}
\end{figure}

\begin{figure}[H]
  \begin{center}
    \includegraphics[width=\textwidth]{/home/urtzi/Documents/papers/individuals_DEB/figures/2_final/figS1-metabolism.pdf}
  \end{center}
  \caption{Model predictions for metabolic rates for the population-average parameterisation (\textit{blue}: males; \textit{red}: females) and eight selected individuals.}\label{fig:fit-metabolism}
\end{figure}

\newpage

\section{Growth trajectories}

Individual 20-year DEB growth simulations at the \textbf{(a)} Czech reference site (\ang{49.2247}\,N, \ang{16.0579}\,E; \qty{600}{m} a.s.l.) and \textbf{(b)} at a comparatively warmer site in Northern Spain (inside the known distribution of the species). A vertical dashed line marks when individuals reached puberty (first hour where maturity \ensuremath{E_H} exceeds the puberty threshold \ensuremath{E_H^p}). We present example simulations for females first (with the population-average parameters, and for four selected individuals) and then for males (population-average parameters first and for individual-level parameters after). Note that the selected individuals correspond to the individuals shown in \textbf{Figs.~\ref{fig:fit-growth}} and \textbf{\ref{fig:fit-metabolism}}.

\newcommand{\growthfigure}[3]{%
  \begin{figure}[H]
    \centering
    \includegraphics[width=\textwidth]{/home/urtzi/Documents/papers/individuals_DEB/figures/1_base/growth_#1.pdf}
    \caption{Growth trajectory for #2.}
    \label{#3}
  \end{figure}%
}

\growthfigure{sp}{the population-average parameterisation (female)}{fig:growth-1}

\growthfigure{8045}{PID 8045 (female)}{fig:growth-2}

\growthfigure{8050}{PID 8050 (female)}{fig:growth-3}

\growthfigure{8069}{PID 8069 (female)}{fig:growth-4}

\growthfigure{8078}{PID 8078 (female)}{fig:growth-5}

\growthfigure{sp_male}{the population-average parameterisation (male)}{fig:growth-6}

\growthfigure{8046}{PID 8046 (male)}{fig:growth-7}

\growthfigure{8053}{PID 8053 (male)}{fig:growth-8}

\growthfigure{8058}{PID 8058 (male)}{fig:growth-9}

\growthfigure{8073}{PID 8073 (male)}{fig:growth-10}

% REFERENCES section
\printbibliography[heading=bibnumbered, title={References}]



\end{document}
